\subsection{\colorbox{orange!80}{\parbox{\textwidth}{PT03-BLD-SQL-INJ: Wydobycie hasha hasła wybranego użytkownika z użyciem Blind SQL Injection}}}
\subsubsection{Opis}
Celem ataku był endpoint służący do logowania. Wykorzystywanym bitem informacji było zalogowanie, bądź jego brak w zależności od dodatkowego warunku wstawionego w pole email. Warunek ten sprawdzał, czy znak na danej pozycji hasha hasła jest równy np. "2". Jeśli to była prawda dochodziło do zalogowania, natomiast w przeciwnym wypadku zwracany był błąd 401 Invalid email or password. 

Atak został przeprowadzony z użyciem prostego skryptu, który iteracyjnie sprawdzał kolejne znaki na kolejnych pozycjach, aż do uzyskania pełnego hasha.
\subsubsection{Warunki niezbędne do wykorzystania podatności}
Wykonanie ataku nie wymaga logowania. Atak wykorzystuje podatność SQL Injection pozwalającą na dodanie funkcji substr do żądania. Atakujący musi posiadać adres e-mail konta, którego hash chce wydobyć.
\subsubsection{Szczegóły wykonania}
W tabeli zawarty został przykład sprawdzający, czy pierwszy znak (odpowiada za to drugi argument funkcji substr) hasha jest równy "2".

\begin{table}[H]
    \begin{tabular}{|l|l|l|l|l|l|}
    \hline
    \rowcolor[HTML]{EFEFEF} 
    \multicolumn{1}{|c|}{\cellcolor[HTML]{EFEFEF}\textbf{Metoda}} & \multicolumn{1}{c|}{\cellcolor[HTML]{EFEFEF}\textbf{Końcówka}} & \multicolumn{1}{c|}{\cellcolor[HTML]{EFEFEF}\textbf{Nagłowki}} & \multicolumn{1}{c|}{\cellcolor[HTML]{EFEFEF}\textbf{Ciasteczka}} & \multicolumn{1}{c|}{\cellcolor[HTML]{EFEFEF}\textbf{Dane}}                             & \multicolumn{1}{c|}{\cellcolor[HTML]{EFEFEF}\textbf{Odpowiedź/wynik}} \\ \hline
    POST                                                          & /rest/user/login                                        & -                                  & -                 & \begin{tabular}[c]{@{}l@{}}\{\\ "email":"admin@juice-sh.op' \\ and substr(password,1,1)='2'  -- ", \\ "password":""\\ \}\end{tabular} & 
    \begin{tabular}[c]{@{}l@{}}200 + TOKEN 
    \end{tabular}
        
    \\ \hline
    \end{tabular}
    \end{table}

Przykład zakończył się logowaniem, czyli skrypt zapisywał znak "2" jako poprawny i przechodził do kolejnej pozycji, aż do napotkania pozycji dla której żaden znak z listy nie był właściwy (czyli hash się skończył).
\subsubsection{Skutki}
Wykonanie skryptu pozwoliło na \textbf{pobranie wartości hasha hasła administratora} (admin@juice-sh.op): 

240be518fabd2724ddb6f04eeb1da5967448d7e831c08c8fa822809f74c720a9

\subsubsection{Zalecenia} 
Problemem tutaj jest przede wszystkim sama możliwość wykonania SQL Injection, która umożliwia dostęp do hasha hasła. Po naprawieniu tej podatności, jeśli atak Blind SQL Injection wciąż byłby możliwy to należało by rozważyć:
\begin{itemize}
    \item dodanie losowych małych opóźnień w celu ukrycia czasu przetwarzania
    \item poprawę mechanizmu informowania o błędach. Użytkownik nie musi wiedzieć dokładnie co poszło nie tak - szczególnie jeśli nie zależy to od dostarczanych przez niego danych. Czasami lepiej odpowiedzieć "Spróbuj ponownie później", a szczegóły dla programistów zapisać w specjalnych logach
\end{itemize}